\thispagestyle{empty}
\linespread{1.0}
\newgeometry{top=2.5cm, bottom=2.5cm, left=3cm, right=3cm}

	\begin{titlepage}
		\def\dystans{1cm}
		\begin{center}
			\Large 
\begin{figure}[h]
    \centering
    \includegraphics[height=3.2cm]{prz_pl.png}%
    \hfill
    \includegraphics[height=3.2cm]{wmifs_pl.png}
\end{figure} 
			{\large\scshape Politechnika Rzeszowska}
            
            {\large im. Ignacego Łukasiewicza}
        
            {\large\scshape\MakeUppercase{Wydział Matematyki i Fizyki Stosowanej}}
		\end{center}
		\vspace{1.3cm}
\begin{center}
  \begin{minipage}[c][5cm][c]{10cm}
    \centering
    {\Large Mikołaj Data} \par\vspace{0.4cm}
    {\Large\textbf{\textsc{Modelowanie dużych sieci LoRa Mesh}}} \par\vspace{0.4cm}
    Promotor: dr Krzysztof Pupka
  \end{minipage}
\end{center}
        
% \begin{center}
%     \begin{minipage}[c][5cm][c]{10cm} % 10 cm szerokości, 5 cm wysokości, wyrównane centralnie
%         \begin{center}
%             {\Large Imię i nazwisko autora}
%         \end{center}    
%         \begin{center}
%             {\Large\textbf{\textsc{Segmentacja zbioru sonogramów stawów biodrowych niemowląt ukierunkowana pod automatyzację oceny morfologii stawu biodrowego}}}
%         \end{center}
%         \begin{center}
%             Promotor: tytuł/stopień naukowy, imię i nazwisko
%         \end{center}
%     \end{minipage}
% \end{center}

\vspace{2cm}

        
		\begin{center}
			{\large  \textsc{Praca  inżynierska} \\
				\vspace{0.5cm}
				\large kierunek studiów:  Inżynieria i analiza danych}
		\end{center}
		\vspace{1.1cm}
		
		

		
		
		\vfill 
		\begin{center}
			{\large Rzeszów 2026}
		\end{center}
	\end{titlepage}
    \restoregeometry